% =====================================================================
%  figures_pgfplots.tex  --  camera-ready vector figures for the
%  privacy-accounting section.  Every plot reads a .dat file written by
%  `python -m sdpa.figures`, so text, ticks and fonts are LaTeX-native.
%
%  Set the data directory before \input, e.g.
%      \def\figdatadir{../figures/dat/}
%  The default below assumes you compile from inside figures/.
% =====================================================================
\providecommand{\figdatadir}{dat/}
\pgfplotsset{compat=1.18}

% ---------------------------------------------------------------------
% Fig. 1 -- the accounting law dictates the noise profile
% ---------------------------------------------------------------------
\begin{figure}[t]
\centering
\begin{tikzpicture}
  \pgfplotstableread{\figdatadir budget_split.dat}\bs
  \begin{axis}[
      width=\columnwidth, height=0.70\columnwidth,
      xlabel={semantic component $h$},
      ylabel={share of noise power},
      xmin=-0.8, xmax=14.4, ymin=0, ymax=0.30,
      xtick={0,2,...,14},
      tick label style={font=\footnotesize},
      label style={font=\footnotesize},
      legend style={font=\scriptsize, draw=none, fill=none,
                    at={(0.5,-0.26)}, anchor=north, legend columns=3},
      ymajorgrids, grid style={densely dotted, black!25},
      extra x ticks={0,1,3,4,5,6},
      extra x tick labels={},
      extra tick style={grid=major, major grid style={black!15}},
    ]
    \addplot[draw=none, fill=gray!10] coordinates {
      (1.5,0)(2.5,0)(2.5,0.30)(1.5,0.30)(1.5,0)
      (6.5,0)(7.5,0)(7.5,0.30)(6.5,0.30)(6.5,0)
      (7.5,0)(8.5,0)(8.5,0.30)(7.5,0.30)(7.5,0)
      (8.5,0)(9.5,0)(9.5,0.30)(8.5,0.30)(8.5,0)
      (9.5,0)(10.5,0)(10.5,0.30)(9.5,0.30)(9.5,0)
      (10.5,0)(11.5,0)(11.5,0.30)(10.5,0.30)(10.5,0)
      (11.5,0)(12.5,0)(12.5,0.30)(11.5,0.30)(11.5,0)
      (12.5,0)(13.5,0)(13.5,0.30)(12.5,0.30)(12.5,0)
      (13.5,0)(14.5,0)(14.5,0.30)(13.5,0.30)(13.5,0)};
    \addplot+[ybar, bar width=4pt, fill=blue!55!black, draw=none]
      table[x expr=\thisrow{h}-0.27, y=gdp] {\bs};
    \addplot+[ybar, bar width=4pt, fill=blue!30, draw=none]
      table[x expr=\thisrow{h}-0.09, y=rdp] {\bs};
    \addplot+[ybar, bar width=4pt, fill=orange!70!black, draw=none]
      table[x expr=\thisrow{h}+0.09, y=basic] {\bs};
    \addplot+[ybar, bar width=4pt, fill=black!40, draw=none]
      table[x expr=\thisrow{h}+0.27, y=static] {\bs};
    \addplot[black, mark=*, mark size=1pt, densely dotted, mark options={solid}]
      table[x=h, y=wkappa_norm] {\bs};
    \legend{GDP, RDP, $\ell_1$ (ampl.), static ($k{=}4$), $w_h\kappa_h$}
  \end{axis}
\end{tikzpicture}
\caption{Noise-power share of the four accounting rules on the task-decoupled
calibration ($\varepsilon{=}100$, $\delta{=}10^{-5}$; this panel illustrates
the shape of the allocation, while the audit in Fig.~\ref{fig:composite}(b)
uses the $\varepsilon{=}2000$ operating point). Only the declared private
components $h\in P=\{0,1,3,4,5,6\}$ may carry noise
(Theorem~\ref{th:allocation}); the gray-shaded blocks are the task complement
$Q$. The amplitude-only $\ell_1$ rule concentrates noise on a few high-sensitivity
components, the quadrature GDP/RDP rules spread it more evenly, and the static
rule adds noise to only the four selected blocks and leaves the rest untouched
(Corollary~\ref{co:static}).}
\label{fig:budget_split}
\end{figure}

% ---------------------------------------------------------------------
% Fig. 2 -- the price of a certificate
% ---------------------------------------------------------------------
\begin{figure}[t]
\centering
\begin{tikzpicture}
  \pgfplotstableread{\figdatadir certificate_price.dat}\cp
  \begin{axis}[
      name=left, width=0.84\columnwidth, height=0.50\columnwidth,
      xlabel={certified error floor $\tau$},
      ylabel={Bob LPIPS},
      xmin=0.29, xmax=0.96, ymin=0.10, ymax=1.0,
      tick label style={font=\footnotesize},
      label style={font=\footnotesize},
      legend style={font=\scriptsize, draw=none, fill=none,
                    at={(0.5,-0.24)}, anchor=north, legend columns=2,
                    nodes={font=\scriptsize, inner sep=1pt}},
      ymajorgrids, grid style={densely dotted, black!25},
    ]
    \addplot[red!70!black, mark=square*, mark size=1.2pt]
      table[x=tau, y=lpips_aligned] {\cp};
    \addplot[blue!70!black, mark=*, mark size=1.2pt]
      table[x=tau, y=lpips_decoupled] {\cp};
    \addplot[black, densely dotted, domain=0.29:0.96] {0.112};
    \legend{{Bob, aligned $w_h$}, {Bob, decoupled $w_h$},
            {no-protection $0.112$}}
  \end{axis}
  \begin{axis}[
      name=right, at={(left.south west)}, anchor=south west,
      width=0.84\columnwidth, height=0.50\columnwidth,
      axis y line*=right, axis x line=none, xtick=\empty,
      ylabel={identification error $P_e$},
      ymin=0.25, ymax=1.15,
      tick label style={font=\footnotesize},
      label style={font=\footnotesize},
    ]
    \addplot[blue!70!black, densely dotted, very thick]
      table[x=tau, y=certified_floor] {\cp};
    \addplot[blue!70!black, only marks, mark=triangle*, mark size=1.4pt,
             densely dashed]
      table[x=tau, y=p_e_surrogate] {\cp};
  \end{axis}
\end{tikzpicture}
\caption{Certificate price on the synthetic anchor, two overlap regimes. The
two Bob LPIPS curves are read on the \emph{left} axis; the dotted
certified-floor curve and the triangle surrogate $P_e$ are read on the
\emph{right} axis. The LPIPS values are surrogate predictions from the default
distortion model, not measured reconstruction distortion. When the private
semantics \emph{are} the task (aligned, $w_h$ large, task overlap $\approx1$)
the certified floor drives Bob's LPIPS to $0.63$--$0.87$, leaving little usable
signal for the legitimate receiver; when the private subspace is task-null
(decoupled) the same floor costs almost nothing until
$\tau$ approaches the ceiling $\tau_{\max}(10^{4}){=}0.925$. The dotted line
is the certified floor itself, which the surrogate $P_e$ must dominate
(Theorem~\ref{th:certificate}).}
\label{fig:certificate}
\end{figure}

% ---------------------------------------------------------------------
% Fig. 4 -- leakage budget I_max(tau, K) sensitivity to the gallery size
% ---------------------------------------------------------------------
\begin{figure}[t]
\centering
\begin{tikzpicture}
  \pgfplotstableread{\figdatadir imax_gallery.dat}\ig
  \begin{axis}[
      width=\columnwidth, height=0.62\columnwidth,
      xlabel={certified error floor $\tau$},
      ylabel={$I_{\max}$ (bit)},
      xmin=0.30, xmax=0.92, ymin=0,
      tick label style={font=\footnotesize},
      label style={font=\footnotesize},
      legend style={font=\scriptsize, draw=none, fill=none,
                    at={(0.98,0.98)}, anchor=north east},
      ymajorgrids, grid style={densely dotted, black!25},
    ]
    \addplot[blue!55!black, dashed] table[x=tau, y=k1e3] {\ig};
    \addplot[blue!70!black, mark=*, mark size=1.2pt] table[x=tau, y=k1e4] {\ig};
    \addplot[orange!80!black, dotted, very thick] table[x=tau, y=k1e6] {\ig};
    \legend{$K=10^{3}$, $K=10^{4}$, $K=10^{6}$}
  \end{axis}
\end{tikzpicture}
\caption{Leakage budget $I_{\max}(\tau,K)$ of \eqref{eq:imax} versus the
certified floor, for three identity-gallery sizes. Raising $K$ pushes the
feasibility ceiling $\tau_{\max}(K)$ and the budget upward only
logarithmically: adjacent gallery sizes differ by about $0.025$ in
$\tau_{\max}$ and the full span over three orders of magnitude ($10^{3}$ to
$10^{6}$) is about $0.050$. Curves are truncated at their own feasibility
ceiling $\tau_{\max}(K)$; beyond it no non-negative budget exists.}
\label{fig:imax}
\end{figure}

% ---------------------------------------------------------------------
% Fig. 3 -- composite, one column per story
% ---------------------------------------------------------------------
\begin{figure*}[t]
\centering
\begin{tikzpicture}
  \pgfplotstableread{\figdatadir waterfill_sweep.dat}\wf
  \begin{axis}[
      width=0.315\textwidth, height=0.27\textwidth,
      xlabel={dual variable $\nu$}, ylabel={$\sigma_h^{2}/g_h$},
      xmode=log, ymode=log, ymin=1e-4, ymax=1e8,
      tick label style={font=\footnotesize},
      label style={font=\footnotesize},
      title={(a) water-filling}, title style={font=\footnotesize},
      legend style={font=\scriptsize, draw=none, fill=none,
                    at={(0.98,0.03)}, anchor=south east},
    ]
    \addplot[blue!55!black] table[x=mu, y=w_1] {\wf};
    \addplot[blue!55!black, densely dashed] table[x=mu, y=w_8] {\wf};
    \addplot[orange!80!black, densely dotted, very thick]
      table[x=mu, y=w_125] {\wf};
    \legend{$w_h$, $w_h/8$, $w_h/125$}
  \end{axis}
\end{tikzpicture}
\hfill
\begin{tikzpicture}
  \pgfplotstableread{\figdatadir accounting_audit.dat}\ac
  \begin{axis}[
      width=0.315\textwidth, height=0.27\textwidth,
      xlabel={noise-to-signal ratio (dB)}, ylabel={certified $\epsilon$},
      ymode=log,
      tick label style={font=\footnotesize},
      label style={font=\footnotesize},
      title={(b) audited $\epsilon$ at one release},
      title style={font=\footnotesize},
      legend style={font=\scriptsize, draw=none, fill=none,
                    at={(0.98,0.03)}, anchor=south east},
    ]
    \addplot[draw=none, fill=red!7] coordinates {
      (9.7,1)(24,1)(24,5000)(9.7,5000)(9.7,1)};
    \node[font=\scriptsize, red!60!black, anchor=north west]
      at (axis cs:10.2,3000) {hazard: $\mu_h\gg c$};
    \addplot[blue!55!black, mark=triangle*, mark size=1.5pt]
      table[x=nsr_db, y=eps_quadrature] {\ac};
    \addplot[green!40!black, mark=square*, mark size=1.2pt]
      table[x=nsr_db, y=eps_l1_exact] {\ac};
    \addplot[orange!80!black, mark=*, mark size=1.2pt, densely dashed]
      table[x=nsr_db, y=eps_l1_stylized] {\ac};
    \legend{quadrature, $\ell_1$ exact, $\ell_1$ ampl.}
  \end{axis}
\end{tikzpicture}
\hfill
\begin{tikzpicture}
  \pgfplotstableread{\figdatadir certificate_gap.dat}\cg
  \begin{axis}[
      width=0.315\textwidth, height=0.27\textwidth,
      xlabel={certified floor $\tau$}, ylabel={$P_e^{\rm sim}-\tau$},
      ymin=-0.05, ymax=0.62,
      tick label style={font=\footnotesize},
      label style={font=\footnotesize},
      title={(c) certificate tightness},
      title style={font=\footnotesize},
    ]
    \addplot[purple!60!black, mark=*, mark size=1.2pt]
      table[x=tau, y=gap] {\cg};
    \addplot[black, densely dotted, domain=0.42:0.93] {0};
  \end{axis}
\end{tikzpicture}
\caption{(a) Water-filling charge: one latent direction, three task weights.
Privacy is cheaper for a task-null direction because $a_h=w_h\kappa_h$
enters the KKT solution \eqref{eq:kkt}. (b) Audit of the accounting laws on
the \emph{same} release at the low-noise (large-$\mu$) operating point certified at
$\varepsilon{=}2000$, $\delta{=}10^{-5}$: quadrature (GDP) sits between the
rigorous tail bound and the amplitude-only $\ell_1$ form, the latter under-
reporting $\varepsilon$ once $\mu_h\gg c$. The red shaded band marks this
large-$\mu$ hazard region, where a Laplace-style $\ell_1$ composition yields an
assessment tighter than the guarantee. (c) The surrogate $P_e$ dominates the
certified floor everywhere: the certificate is conservative, not optimistic.}
\label{fig:composite}
\end{figure*}

% ---------------------------------------------------------------------
% Fig. 4 -- latent-level measurement on the released SemanticStyleGAN
% ---------------------------------------------------------------------
\begin{figure*}[t]
\centering
\begin{tikzpicture}
  \pgfplotstableread{\figdatadir latent_utility_landscape.dat}\lu
  \begin{axis}[
      width=0.315\textwidth, height=0.27\textwidth,
      ybar, bar width=3.2pt,
      ylabel={measured $w_h$},
      ymin=0, ymax=0.40,
      xtick={0,...,14},
      xticklabels={sh0,sh1,bg,skin,eye,brw,mth,nos,ear,hair,nck,clth,eyg,hat,erR},
      x tick label style={rotate=55, anchor=east, font=\tiny},
      tick label style={font=\footnotesize},
      label style={font=\footnotesize},
      title={(a) measured utility landscape}, title style={font=\footnotesize},
    ]
    \addplot[fill=blue!45, draw=blue!55!black] table[x=h, y=w_h] {\lu};
  \end{axis}
\end{tikzpicture}
\hfill
\begin{tikzpicture}
  \pgfplotstableread{\figdatadir latent_sweep_examples.dat}\ls
  \begin{axis}[
      width=0.315\textwidth, height=0.27\textwidth,
      xlabel={injected $\sigma_h^{2}$}, ylabel={Bob LPIPS},
      xmode=log,
      tick label style={font=\footnotesize},
      label style={font=\footnotesize},
      title={(b) per-block noise sweeps}, title style={font=\footnotesize},
      legend style={font=\scriptsize, draw=none, fill=none,
                    at={(0.03,0.97)}, anchor=north west},
    ]
    \addplot[blue!55!black, mark=*, mark size=1.2pt]
      table[x=sigma2, y=shared0] {\ls};
    \addplot[orange!80!black, mark=square*, mark size=1.2pt]
      table[x=sigma2, y=skin] {\ls};
    \addplot[green!40!black, mark=triangle*, mark size=1.5pt]
      table[x=sigma2, y=hair] {\ls};
    \addplot[purple!60!black, mark=diamond*, mark size=1.5pt, densely dashed]
      table[x=sigma2, y=eyes] {\ls};
    \legend{shared\_0, skin, hair, eyes}
  \end{axis}
\end{tikzpicture}
\hfill
\begin{tikzpicture}
  \pgfplotstableread{\figdatadir latent_kappa.dat}\lk
  \begin{axis}[
      width=0.315\textwidth, height=0.27\textwidth,
      xlabel={injected $\sigma_h^{2}$}, ylabel={residual fraction $\kappa_h$},
      xmode=log, ymin=0, ymax=1.05,
      tick label style={font=\footnotesize},
      label style={font=\footnotesize},
      title={(c) Wiener deprotector residual}, title style={font=\footnotesize},
      legend style={font=\scriptsize, draw=none, fill=none,
                    at={(0.97,0.97)}, anchor=north east},
    ]
    \addplot[blue!55!black, mark=*, mark size=1.2pt]
      table[x=sigma2, y=kappa_empirical] {\lk};
    \addplot[red!70!black, densely dashed, very thick]
      table[x=sigma2, y=kappa_analytic] {\lk};
    \legend{empirical (held-out), analytic}
  \end{axis}
\end{tikzpicture}
\caption{Latent-level measurement on the released SemanticStyleGAN generator.
(a) Measured per-block utility weights $w_h$: the perceptual cost of latent
noise concentrates on the shared codes, hair, skin and background, while the
facial-feature blocks (eyes, nose, ear\_r) are nearly free to protect. Block
order: shared\_0, shared\_1, background, skin, eyes, brows, mouth, nose, ears,
hair, neck, cloth, eye\_g, hat, ear\_r. (b) Example decode-and-LPIPS sweeps
with the quadratic fit $\mathrm{LPIPS}\approx\mathrm{lpips}_0+w_h\sigma_h^2$;
the eyes curve is almost flat---the generator absorbs local latent noise into
plausible variations. (c) Residual fraction of the per-block Wiener
deprotector on held-out codes versus the analytic prediction; the two agree to
three decimals, so $\kappa_h$ is a measured quantity, not a fit artefact.}
\label{fig:latent}
\end{figure*}
